%% toc: P1 4 4.5,4.6
\section{Finite tuples and finite sequences over an infinite cardinal}
\label{sec:p01-w02-p01-011-finite-coding}

The theorem on finite coding over an infinite cardinal
(\cref{thm:infinite-cardinal-finite-coding}) and the syntax-counting
corollary (\cref{cor:syntax-counting}) carry all the cardinal arithmetic that
the Henkin construction of \cref{ch:fol-metatheory} and the
L\"owenheim--Skolem theorems (\cref{thm:downward-ls}) consume. Where they are
stated, the theorem carries only a citation. This section proves every clause
of the theorem, and \cref{sec:p01-w02-p01-011-syntax-counting} proves the
corollary from it.

\emph{Inputs.} The definitions of \cref{sec:set-prerequisites-05}
(finite and countable sets, cardinal comparison by injections and
bijections), the Cantor--Schr\"oder--Bernstein theorem and Cantor's theorem
proved there, transitive sets, ordinals and cardinals
(\cref{def:transitive-set,def:ordinal-cardinal}), the sets \(A^n\) of
\cref{sec:set-prerequisites-02}, and, for the clauses that need it, the
Axiom of Choice in the form defined in \cref{sec:set-prerequisites-06}.
The ambient theory is the classical set theory of the metalanguage fixed at
the start of this part; separation and replacement are used to form the
sets named below. Natural numbers are the sets
\(n=\{0,\ldots,n-1\}\) of \cref{sec:set-prerequisites-02}: \(0=\varnothing\),
\(n+1=n\cup\{n\}\), and \(\N\) is the least set containing \(0\) and closed
under \(n\mapsto n+1\), so that statements about all \(n\in\N\) are proved by
induction on \(n\).

\emph{Conclusions.} Elementary facts about ordinals and finite sets
(\cref{thm:p01-w02-p01-011-ordinals,thm:p01-w02-p01-011-finite-sets,thm:p01-w02-p01-011-infinite-ordinals}),
the representation of initial segments of a well-order by ordinals
(\cref{thm:p01-w02-p01-011-representation}), the identity
\(\abs{\mu\times\mu}=\abs\mu\) for every infinite cardinal \(\mu\) without
any use of Choice (\cref{thm:p01-w02-p01-011-square}), and from these the
four clauses of the finite-coding theorem
(\cref{thm:p01-w02-p01-011-finite-coding}).

\subsection{Readings and order vocabulary}
\label{sec:p01-w02-p01-011-readings}

\begin{hbconvention}{def:p01-w02-p01-011-readings}{Cardinal readings and ordinal vocabulary}
\begin{enumerate}
  \item For sets \(A\) and \(B\), \(\abs A\leq\abs B\), \(\abs A=\abs B\) and
        \(\abs A<\abs B\) have the meanings of \cref{sec:set-prerequisites-05}:
        an injection \(A\to B\) exists; a bijection \(A\to B\) exists; the
        first holds and the second fails. When \(B\) is an ordinal
        \(\kappa\), these are written \(\abs A\leq\kappa\),
        \(\abs A=\kappa\) and \(\abs A<\kappa\). Under the Axiom of Choice,
        \(\abs A\) also names the cardinal bijective with \(A\)
        (\cref{cor:cardinal-assignment}); item~7 of
        \cref{thm:p01-w02-p01-011-infinite-ordinals} shows that for a
        cardinal \(\kappa\) the two readings of \(\abs A<\kappa\) agree, so
        each clause below may be read either way.
  \item For a set \(A\) and \(n\in\N\), \(A^n\) is the set of functions
        \(n\to A\). The set of finite sequences from \(A\) is
        \(A^{<\omega}=\bigcup_{n\in\N}A^n\); it is the set \(A^*\) of strings
        over \(A\) of \cref{sec:strings-01}. Two finite sequences are equal
        iff they have the same domain \(n\) and the same value at every
        \(i<n\). For \(s\in A^{n+1}\), \(s\mathbin{\upharpoonright}n\) is the
        restriction of \(s\) to \(n\).
  \item For ordinals \(\alpha,\beta\): \(\alpha<\beta\) iff
        \(\alpha\in\beta\) (\cref{def:ordinal-cardinal});
        \(\alpha\leq\beta\) iff \(\alpha<\beta\) or \(\alpha=\beta\);
        \(\alpha+1=\alpha\cup\{\alpha\}\); and \(\max(\alpha,\beta)\) is
        \(\beta\) if \(\alpha\leq\beta\) and \(\alpha\) otherwise. An
        ordinal is \emph{infinite} iff it is bijective with no \(n\in\N\);
        for cardinals this is the definition in \cref{def:ordinal-cardinal}.
  \item Let \(\prec\) strictly well-order a set \(W\) in the sense of
        \cref{sec:set-prerequisites-03}. For \(w\in W\), the \emph{initial
        segment} below \(w\) is
        \(\operatorname{seg}(w)=\{v\in W:v\prec w\}\), ordered by the
        restriction of \(\prec\). An \emph{order isomorphism} from
        \((W,\prec)\) to a strictly ordered set \((V,<)\) is a bijection
        \(f:W\to V\) with \(v\prec w\) iff \(f(v)<f(w)\); an ordinal is
        ordered by \(\in\).
\end{enumerate}
\end{hbconvention}

\subsection{Ordinals, finite sets and infinite ordinals}
\label{sec:p01-w02-p01-011-basic}

The three lemmas of this subsection supply the facts about ordinals and
finite sets that the counting arguments use. The inherited text states
ordinal trichotomy later, in the set-theory part; here it is proved from
\cref{def:ordinal-cardinal} alone.

\begin{hblemma}{thm:p01-w02-p01-011-ordinals}{proved=annex:p01-w02-p01-011-ordinals}{Basic facts about ordinals}
Let \(\alpha\) and \(\beta\) be ordinals.
\begin{enumerate}
  \item Every member of \(\alpha\) is an ordinal, and \(\alpha\notin\alpha\).
  \item If \(\alpha\subseteq\beta\), then \(\alpha=\beta\) or
        \(\alpha\in\beta\).
  \item Exactly one of \(\alpha\in\beta\), \(\alpha=\beta\),
        \(\beta\in\alpha\) holds, and \(\alpha\leq\beta\) iff
        \(\alpha\subseteq\beta\).
  \item Every nonempty set \(S\) of ordinals has a unique least member: a
        \(\gamma\in S\) with \(\gamma\leq\delta\) for every \(\delta\in S\).
  \item Every transitive set of ordinals is an ordinal. In particular
        \(\alpha+1\) is an ordinal, and an ordinal \(\gamma\) satisfies
        \(\gamma<\alpha+1\) iff \(\gamma\leq\alpha\).
  \item Every \(n\in\N\) is an ordinal; hence \(n\notin n\), and for
        \(m,n\in\N\), \(m+1=n+1\) implies \(m=n\).
\end{enumerate}
\end{hblemma}

\begin{hblemma}{thm:p01-w02-p01-011-finite-sets}{proved=annex:p01-w02-p01-011-finite-sets}{Finite sets}
\begin{enumerate}
  \item If \(\abs C=\abs r\) with \(r\in\N\) and \(x\notin C\), then
        \(\abs{C\cup\{x\}}=\abs{r+1}\).
  \item The union of two disjoint finite sets is finite.
  \item Every subset of a finite set is finite.
  \item If \(k\in\N\) and \(A_i\) is finite for every \(i<k\), then
        \(\bigcup_{i<k}A_i\) is finite.
  \item If \(A\) and \(B\) are finite, then \(A\times B\) is finite.
\end{enumerate}
\end{hblemma}

\begin{hblemma}{thm:p01-w02-p01-011-infinite-ordinals}{proved=annex:p01-w02-p01-011-infinite-ordinals}{Infinite ordinals and cardinals}
Let \(\alpha\) be an infinite ordinal and \(\kappa\) an infinite cardinal.
\begin{enumerate}
  \item \(\N\subseteq\alpha\). Consequently every countable set \(X\)
        satisfies \(\abs X\leq\alpha\); in particular
        \(\aleph_0=\abs\N\leq\kappa\).
  \item \(\abs{\alpha+1}=\abs\alpha\).
  \item If \(\gamma<\kappa\), then \(\gamma+1<\kappa\).
  \item For every ordinal \(\beta\) there is a least ordinal
        \(\operatorname{card}(\beta)\) bijective with \(\beta\). It satisfies
        \(\operatorname{card}(\beta)\leq\beta\), it is a cardinal, and it is infinite iff
        \(\beta\) is infinite.
  \item If \(\beta<\kappa\), then every set \(X\) with
        \(\abs X\leq\abs\beta\) satisfies \(\abs X<\kappa\).
  \item Every finite set \(X\) satisfies \(\abs X<\kappa\).
  \item If \(\mu\) is a cardinal and \(\abs X=\abs\mu\), then
        \(\abs X<\kappa\) iff \(\mu<\kappa\).
\end{enumerate}
\end{hblemma}

Item~1 is the fact \(\aleph_0\leq\kappa\) that the syntax-counting
corollary needs when it adjoins at most countably many symbols to an
alphabet. Item~4 does not use the Axiom of Choice: an ordinal is already
well-ordered, so its cardinal is found by taking a least member of a set of
ordinals.

\subsection{Well-orders and the square of an infinite cardinal}
\label{sec:p01-w02-p01-011-square}

\begin{hbdefinition}{def:p01-w02-p01-011-canonical-order}{Canonical order on a square of ordinals}
Let \(\mu\) be an ordinal. For \((\alpha,\beta),(\gamma,\delta)\in\mu\times\mu\),
put \((\alpha,\beta)\lhd(\gamma,\delta)\) iff one of the following holds:
\begin{enumerate}
  \item \(\max(\alpha,\beta)<\max(\gamma,\delta)\);
  \item \(\max(\alpha,\beta)=\max(\gamma,\delta)\) and \(\alpha<\gamma\);
  \item \(\max(\alpha,\beta)=\max(\gamma,\delta)\), \(\alpha=\gamma\) and
        \(\beta<\delta\).
\end{enumerate}
Equivalently, \(p\lhd q\) iff the triple
\(k(p)=(\max(\alpha,\beta),\alpha,\beta)\) for \(p=(\alpha,\beta)\) precedes
\(k(q)\) in the lexicographic order of triples of ordinals.
\end{hbdefinition}

\begin{hblemma}{thm:p01-w02-p01-011-canonical-order}{proved=annex:p01-w02-p01-011-canonical-order}{The canonical order}
Let \(\mu\) be an ordinal.
\begin{enumerate}
  \item \(\lhd\) strictly well-orders \(\mu\times\mu\).
  \item If \((\xi,\eta)\lhd(\alpha,\beta)\) and
        \(\gamma=\max(\alpha,\beta)\), then \(\xi,\eta\in\gamma+1\). Hence
        \(\operatorname{seg}(\alpha,\beta)\subseteq(\gamma+1)\times(\gamma+1)\).
\end{enumerate}
\end{hblemma}

\begin{hblemma}{thm:p01-w02-p01-011-representation}{proved=annex:p01-w02-p01-011-representation}{Initial segments are represented by ordinals}
Let \(\prec\) strictly well-order a set \(W\).
\begin{enumerate}
  \item If \(f:\alpha\to\beta\) is an order isomorphism between ordinals,
        then \(\alpha=\beta\) and \(f\) is the identity.
  \item For every \(w\in W\) there is a unique ordinal \(F(w)\) such that
        \(\operatorname{seg}(w)\) is order isomorphic to \(F(w)\). If
        \(v\prec w\), then \(F(v)\in F(w)\); so \(F\) is injective.
  \item Let \(\kappa\) be a cardinal. If
        \(\abs{\operatorname{seg}(w)}<\kappa\) for every \(w\in W\), then
        \(F(w)\in\kappa\) for every \(w\in W\), and \(F\) is an injection
        \(W\to\kappa\).
\end{enumerate}
\end{hblemma}

\begin{hbtheorem}{thm:p01-w02-p01-011-square}{proved=annex:p01-w02-p01-011-square}{The square of an infinite cardinal}
For every infinite cardinal \(\mu\), \(\abs{\mu\times\mu}=\abs\mu\). The
proof uses no form of the Axiom of Choice.
\end{hbtheorem}

\begin{hbexample}{def:p01-w02-p01-011-omega-square}{The canonical order on \(\N\times\N\)}
Take \(\mu=\N\). It is an ordinal: every \(n\in\N\) is an ordinal by item~6
of \cref{thm:p01-w02-p01-011-ordinals}, every member of \(n\in\N\) is a
natural number, and a transitive set of ordinals is an ordinal by item~5
there. The pairs with maximum \(\gamma\) form the group
\[
  (0,\gamma),(1,\gamma),\ldots,(\gamma-1,\gamma),(\gamma,0),(\gamma,1),\ldots,(\gamma,\gamma)
\]
of \(2\gamma+1\) pairs, listed in \(\lhd\)-order, and every pair of the group
follows the \(\gamma^2\) pairs whose maximum is smaller. The order therefore
begins
\[
  (0,0),\ (0,1),\ (1,0),\ (1,1),\ (0,2),\ (1,2),\ (2,0),\ (2,1),\ (2,2),\ (0,3),\ \ldots
\]
and the position of a pair, counted from \(0\), is
\[
  F(\alpha,\beta)=
  \begin{cases}
    \beta^2+\alpha,&\alpha<\beta,\\
    \alpha^2+\alpha+\beta,&\beta\leq\alpha.
  \end{cases}
\]
Group \(\gamma\) occupies exactly the positions \(\gamma^2,\ldots,\gamma^2+2\gamma=(\gamma+1)^2-1\),
so consecutive groups tile \(\N\) and \(F\) is a bijection
\(\N\times\N\to\N\); for example \(F(1,1)=3\), \(F(0,2)=4\) and
\(F(2,1)=7\). This is a second explicit pairing beside the map
\(\operatorname{pair}_{\N}\) of \cref{sec:strings-01}, and it is the
function \(F\) of \cref{thm:p01-w02-p01-011-representation} for
\((\N\times\N,\lhd)\): its value at a pair is the number of pairs
preceding it. For an uncountable \(\mu\) no such formula is
available; the proof of \cref{thm:p01-w02-p01-011-square} replaces counting
by the bound on initial segments in
\cref{thm:p01-w02-p01-011-canonical-order}.
\end{hbexample}

\subsection{Finite coding over an infinite cardinal}
\label{sec:p01-w02-p01-011-coding}

The next theorem is \cref{thm:infinite-cardinal-finite-coding} with the same
hypotheses and clauses, stated again so that its status points to the local
proof.

\begin{hbtheorem}{thm:p01-w02-p01-011-finite-coding}{proved=annex:p01-w02-p01-011-finite-coding}{Finite coding over an infinite cardinal}
Assume the Axiom of Choice, and let \(\kappa\) be an infinite cardinal.
Then:
\begin{enumerate}
  \item \(\abs{\kappa^n}=\kappa\) for every positive \(n\in\N\);
  \item the set of all finite sequences from \(\kappa\) has cardinality
        \(\kappa\);
  \item if \(\alpha<\kappa\) and \(A_\beta\) is finite for every
        \(\beta<\alpha\), then
        \(\abs{\bigcup_{\beta<\alpha}A_\beta}<\kappa\); and
  \item if \(\abs{A_n}\leq\kappa\) for every \(n\in\N\), then
        \(\abs{\bigcup_{n\in\N}A_n}\leq\kappa\).
\end{enumerate}
\end{hbtheorem}

\emph{Where Choice enters the proof.} Clause~1 rests on
\cref{thm:p01-w02-p01-011-square} and uses no choice. The proof of clause~2
uses the Axiom of Choice once, to select one injection
\(\kappa^n\to\kappa\) for every positive \(n\). Clause~3 for a finite
\(\alpha\) uses no choice; for an infinite \(\alpha\) the proof selects an
enumeration of every \(A_\beta\). Clause~4 is the one place where the
hypothesis gives, for each \(n\), only the existence of an injection
\(A_n\to\kappa\); the proof applies the Axiom of Choice to the family of the
sets of such injections to obtain the whole sequence
\((f_n)_{n\in\N}\) at once. These remarks describe what the proofs below use;
they are not claims that the clauses fail without Choice.

\begin{hbwarning}{def:p01-w02-p01-011-coding-limits}{The hypotheses of the finite-coding theorem are used}
\begin{enumerate}
  \item \emph{Clause~3 needs \(\alpha<\kappa\).} For \(\alpha=\kappa\) and
        \(A_\beta=\{\beta\}\), every \(A_\beta\) is finite but
        \(\bigcup_{\beta<\kappa}A_\beta=\kappa\), which is not of
        cardinality \({<}\kappa\).
  \item \emph{Clause~3 needs finite pieces.} For \(\alpha=1\) and
        \(A_0=\kappa\), the union is \(\kappa\). Clause~4 gives only
        \(\leq\kappa\) for a countable union of sets of size
        \(\leq\kappa\), and the bound is attained: for \(A_n=\kappa\) the
        union is \(\kappa\).
  \item \emph{Clause~2 needs finite length.} The set of all functions
        \(\N\to\kappa\) can be larger than \(\kappa\). For
        \(\kappa=\aleph_0\), the cardinal of \(\N\)
        (\cref{sec:set-prerequisites-06}), consider the functions
        \(\N\to\N\). Sending \(S\subseteq\N\) to its characteristic function
        \(\N\to\{0,1\}\subseteq\N\) injects \(\mathcal P(\N)\) into them; an
        injection of them into \(\N\) would then give an injection
        \(g:\mathcal P(\N)\to\N\), and the map sending \(n\) to the unique \(S\)
        with \(g(S)=n\) when there is one, and to \(\varnothing\) otherwise,
        would be a surjection \(\N\to\mathcal P(\N)\), which Cantor's theorem in
        \cref{sec:set-prerequisites-05} excludes.
\end{enumerate}
\end{hbwarning}

\section{Counting expressions in arbitrary set-sized languages}
\label{sec:p01-w02-p01-011-syntax-counting}

A language whose nonlogical vocabulary is a set of size \(\kappa\) has, by
the corollary below, at most \(\kappa\) strings, tuples of strings and
formulas when \(\kappa\) is infinite. The proof needs two facts that are easy
to pass over: an expression produced by a grammar whose rules are finitary is
a finite string over the grammar's alphabet, and adjoining countably many
logical symbols does not raise the bound. The first is
\cref{thm:p01-w02-p01-011-grammar-strings}; the second is item~1 of
\cref{thm:p01-w02-p01-011-infinite-ordinals}.

\begin{hbdefinition}{def:p01-w02-p01-011-grammar}{Finitary grammar over an alphabet}
Let \(\Gamma\) be a set, the \emph{alphabet}. A \emph{finitary grammar over
\(\Gamma\)} is a pair \((B,\mathcal F)\) of a set \(B\subseteq\Gamma^*\) of
\emph{base strings} and a set \(\mathcal F\) of \emph{formation rules},
each \(f\in\mathcal F\) a function \(f:(\Gamma^*)^{k_f}\to\Gamma^*\) for some
\(k_f\in\N\), its arity. The set of \emph{expressions generated} by the
grammar is the intersection of all sets \(Y\subseteq\Gamma^*\) with
\(B\subseteq Y\) that are closed under every rule in \(\mathcal F\). This is
the inductively generated set of \cref{sec:strings-02} with universe
\(U=\Gamma^*\), base set \(B\) and operations \(\mathcal F\). A rule given by
a \emph{template} \(u_0,\ldots,u_k\in\Gamma^*\) is the function
\[
  (x_1,\ldots,x_k)\mapsto u_0x_1u_1x_2\cdots x_ku_k ,
\]
where juxtaposition is the concatenation of \cref{sec:strings-01},
associated to the left.
\end{hbdefinition}

\begin{hblemma}{thm:p01-w02-p01-011-grammar-strings}{proved=annex:p01-w02-p01-011-grammar-strings}{Generated expressions are finite strings}
Let \(\Gamma\) be an alphabet.
\begin{enumerate}
  \item If \(u\in\Gamma^m\) and \(v\in\Gamma^n\), then
        \(uv\in\Gamma^{m+n}\). Consequently every template rule is a
        function \((\Gamma^*)^k\to\Gamma^*\), hence a formation rule in the
        sense of \cref{def:p01-w02-p01-011-grammar}.
  \item Every expression generated by a finitary grammar over \(\Gamma\)
        is a finite string over \(\Gamma\): the generated set \(G\)
        satisfies \(G\subseteq\Gamma^*=\Gamma^{<\omega}\).
  \item If \(\Gamma\subseteq\Gamma'\), then
        \(\Gamma^*\subseteq\Gamma'^*\).
\end{enumerate}
\end{hblemma}

The next corollary is \cref{cor:syntax-counting} with the same hypotheses
and items; its proof below follows the inherited proof and writes out the
two steps that the inherited proof left to other places.

\begin{hbcorollary}{thm:p01-w02-p01-011-syntax-counting}{proved=annex:p01-w02-p01-011-syntax-counting}{Syntax counting}
Assume the Axiom of Choice and let \(\kappa\) be an infinite cardinal.
\begin{enumerate}
  \item Every set \(A\) with \(\abs{A}\leq\kappa\) satisfies
        \(\abs{A^{<\omega}}\leq\kappa\), where \(A^{<\omega}\) is the set of
        finite sequences from \(A\).
  \item If \(\Sigma\) is an alphabet with \(\abs{\Sigma}\leq\kappa\), then
        the finite strings over \(\Sigma\), the finite tuples of such
        strings, and the formulas generated by any fixed finitary grammar
        whose alphabet extends \(\Sigma\) by at most countably many further
        symbols each form a set of cardinality at most \(\kappa\).
\end{enumerate}
\end{hbcorollary}

\begin{hbexample}{def:p01-w02-p01-011-constants-example}{A language with \(\kappa\) constants}
Let \(\kappa\) be an infinite cardinal and let
\(\Sigma=\{c_\xi:\xi<\kappa\}\) be a set of distinct constant symbols, so
\(\abs\Sigma=\kappa\) by \(\xi\mapsto c_\xi\). Adjoin the set
\(X=\{\doteq,\neg,\wedge,(,)\}\) of five further symbols, none in
\(\Sigma\), and let \(\Gamma=\Sigma\cup X\). Take as base strings the
three-symbol strings \(c_\xi\doteq c_\eta\) for \(\xi,\eta<\kappa\), and as
formation rules the templates \(x\mapsto\neg x\) and
\((x,y)\mapsto(x\wedge y)\), that is, \(u_0=\neg\), \(u_1=\epsilon\) and
\(u_0=(\), \(u_1=\wedge\), \(u_2=)\). This is an illustrative grammar, not a
logic defined elsewhere in the handbook. Item~2 of
\cref{thm:p01-w02-p01-011-syntax-counting} gives at most \(\kappa\)
generated expressions. Conversely \(\xi\mapsto c_\xi\doteq c_\xi\) is an
injection \(\kappa\to G\), because two strings that agree at position
\(0\) have the same first symbol. The Cantor--Schr\"oder--Bernstein theorem
gives \(\abs G=\kappa\): the bound of the corollary is attained. The same
computation with \(X\) countably infinite, for instance with variables
\(v_n\) for \(n\in\N\), changes nothing, because item~1 of
\cref{thm:p01-w02-p01-011-infinite-ordinals} gives \(\abs X\leq\kappa\).
\end{hbexample}

\begin{hbwarning}{def:p01-w02-p01-011-infinitary-limit}{Finitary rules and finite strings are both used}
The bound depends on every rule taking finitely many arguments and returning
a finite string. A formation step that takes an infinite sequence of
expressions, such as a conjunction of countably many formulas, does not
produce a member of \(\Gamma^*\), and \cref{thm:p01-w02-p01-011-grammar-strings}
does not apply to it. Item~3 of \cref{def:p01-w02-p01-011-coding-limits}
shows what is lost: already the infinite sequences over the two-symbol
alphabet \(\{0,1\}\) cannot be injected into \(\N\). The bound is also only
an upper bound: when \(\abs\Sigma<\kappa\), item~2 still gives
\(\leq\kappa\), and for finite or countable \(\Sigma\) the countable case of
\cref{sec:strings-01} gives the sharper bound \(\aleph_0\).
\end{hbwarning}

\paragraph{Relationships and sources.}
The statements of \cref{thm:p01-w02-p01-011-finite-coding} and
\cref{thm:p01-w02-p01-011-syntax-counting} are those of the frozen reference
\cite[ch.~``Sets, Relations, Functions, and Cardinality'', section ``Choice
principles used later'', theorem ``Finite coding over an infinite cardinal''
and corollary ``Syntax counting'']{FrozenFormalLogicFiniteCoding}, which appear
in this book as \cref{thm:infinite-cardinal-finite-coding} and
\cref{cor:syntax-counting}. Every proof and lemma in this section and in
\cref{sec:p01-w02-p01-011-finite-coding} is an editorial completion built from
the definitions of that chapter and of the chapter on finite strings
\cite[ch.~``Finite Strings, Induction, and Recursion'']{FrozenFormalLogicFiniteCoding};
none is attributed to an external presentation. For \(\kappa=\aleph_0\),
item~1 of the corollary is the countability of strings over a countable
alphabet proved in \cref{sec:strings-01} with the pairing
\(\operatorname{pair}_{\N}\). \Cref{thm:p01-w02-p01-011-square} is the
special case \(\kappa=\lambda\) of the identity
\(\kappa\cdot\lambda=\max(\kappa,\lambda)\) in
\cref{thm:p06-cardinal-arithmetic-laws}, proved here without Choice for
\(\kappa=\lambda\), with \(\kappa\cdot\kappa\) read as
\(\abs{\kappa\times\kappa}\). Clause~4 is the countable-union bound whose
use of Choice is listed in the choice-dependency table of Part~6. The
consumers are the Henkin construction (\cref{thm:fol-completeness}) and
\cref{thm:downward-ls}.
